\scanpage{161}
\mspage{80}
% Lignes 1--5 (continuation de la note marginale de la page précédente).
\illegible{2}, and \unsure{to the} \deleted{\unsure{the}} fact that they remain \inserted{very} close to the language \unsure{and|\deleted{and}} to topological intuition, it is superior to the $p^t$ of view ``groups with loops'', which corresponds to the computational approach rather.)

% Lignes 5--16 : grand passage encadré et barré par des traits obliques.
\deleted{, let us be given \deleted{a subset subtopos} (\inserted{\unsure{necessarily$^t$}} \unsure{open}) $B_{D_1^*}$ of $B_{D^*}$ (corresponding to a subset $I_1$ of $I=\pi_0(B_{D^*})$) so that
\[
 B_{D^*}\simeq B_{D_0^*}\amalg B_{D_1^*}
\]
\margin{($D_0^*$ corresponds to the subset $I_0=I\setminus I_1$ of $I$)}
We have an analogous decomposition
\[
 B_D\simeq B_{D_0}\amalg B_{D_1}.
\]
\inserted{\unsure{definition} $B_{\pi_0(D^*)}=B_S$}
Consider the topos \unsure{$B_U'$} ``amalgamated sum'' of \deleted{\unsure{$B_U$}}
\diagram{\begin{tikzcd}[ampersand replacement=\&]
 \& B_{D_1^*}\arrow[dl]\arrow[dr] \& \& \\
 B_U \& \& B_D\arrow[r,leftrightarrow,"\unsure{\sim}"] \& B_S
\end{tikzcd}}
}

% Lignes 17--24.
I realize that I need to return to the notation for the \deleted{\illegible{1}} \inserted{various} topoi associated with such a datum. But I am going to be guided by the situation of numbers$^\circ$ .-- where we have a regular scheme \inserted{$X$} of dimension 1, a closed subscheme \inserted{$S$} of dimension 0, whence $U=X-S$ --- in this case \deleted{\illegible{2}} $\unsure{S_{\text{ét}}|X_{\text{ét}}}$ can \textcircled{not} be reconstructed from $B_{D^*}$ as $\unsure{B_{\pi_0 B_{D^*}}|B_I}$
\mspage{160}


\scanpage{162}
% Lignes 1--7.
with $I=\pi_0(B_{D^*})$, we must take into account the residue fields $k(s)$ $(s\in S)$ i.e. the Galois groups $\operatorname{Gal}(\overline{k(s)}/k(s))$. So there is reason to return to a starting situation (which is adapted to the \inserted{arithmetico-}geometric case) of morphisms of \uuline{multi-Galois} topoi
\[
 (*)\qquad
 \diagram{\begin{tikzcd}[ampersand replacement=\&]
 \& B_{D^*}\arrow[dl,"\rho"']\arrow[dr,"\sigma"] \& \\
 B_U \& \& B_S
 \end{tikzcd}}
\]
% Lignes 8--14.
(Careful, we write $B_S$, not $B_D$, which will have another meaning), \uline{where} $\sigma$ \uline{induces an isomorphism on the} $\Pi_0$ [and is surjective on the $\operatorname{Hom}$]. For the moment, we are not going to make any particular hypothesis on this situation, which for $B_U$ connected, and in terms of the choices
% Lignes 15--23 : crochet et parenthèses tels que visibles.
[\deleted{of objects $\widetilde D_i^*\in\unsure{D}$} of universal coverings $\widetilde D_i^*$ of the connected components $D_i^*$ of $D^*$, and of a universal covering $\widetilde U$ of $B_U$, and of the isomorphism between the $\rho_!(\widetilde D_i^*)$ and $\widetilde U$) is made explicit by the datum of a group \inserted{$\mathcal E$ (or} $\Pi$) $(=\operatorname{Aut}(\widetilde U))$ \deleted{\unsure{and}} \inserted{$\mathcal E_i$ (or} $\Pi_i$ $(=\operatorname{Aut}(\widetilde D_i^*))$ \deleted{\illegible{2}} $(i\in I)$, $\Gamma_i$ $(i\in I)$
\mspage{161}


\scanpage{163}
\mspage{81}
% Lignes 1--3 et premier diagramme.
and group homomorphisms
\diagram{\begin{tikzcd}[ampersand replacement=\&,row sep=2em,column sep=2.5em]
\& \unsure{\mathcal E_i} \arrow[dl,no head,"\rho_i"'] \arrow[dr,"\sigma_i"] \& \\
\unsure{\mathcal E} \& \& \Gamma_i
\end{tikzcd}}
with the $\sigma_i$ surjective. \deleted{\unsure{and}}
% Lignes 3--10.
[When there is a base field $K$ for the geometric situation, then in the \unsure{topos-theoretic} description, \inserted{\unsure{putting} $e=\operatorname{Spec}K$ and} denoting by $B_e$ the étale topos \inserted{$e_{\text{ét}}$} of $e$, i.e. $B_\Gamma$ \unsure{where|if} $\Gamma=\operatorname{Gal}(\overline K/K)$, the diagram $*$ fits into
\diagram{\begin{tikzcd}[ampersand replacement=\&]
 \& B_{D^*}\arrow[dl,"\rho"']\arrow[dr,"\sigma"] \& \\
 B_U\arrow[dr,"\varphi"'] \& \& B_S\arrow[dl,"\psi"] \\
 \& B_e \&
\end{tikzcd}},
with a commutativity datum for the square considered\unsure{\ldots\ as considered} at the beginning of these notes].

% Lignes 11--18 et dernier diagramme.
In terms of $*$, we construct by ``gluing'' of $B_U$ and \inserted{of} $B_S$ (via the gluing functor $\sigma_*\rho^*$) a mixed topos, which is no longer multi-Galois in general, previously denoted $\unsure{D_{X,U}|B_{X,U}}$, and which \inserted{\unsure{I prefer}} now \deleted{\illegible{1}} \unsure{to denote|denoted} $B_{X;S}$ \deleted{\illegible{2}} \inserted{[to recall that these are sheaves on $X$, but having singularities only on $S$]} \deleted{\illegible{1} fits into a}. It fits into a diagram of topoi
\diagram{\begin{tikzcd}[ampersand replacement=\&,row sep=2em,column sep=2.5em]
\& B_{D^*} \arrow[dl] \arrow[dr] \& \\
B_U \arrow[dr,""{name=xl,description}] \& \& B_S \arrow[dl,""{name=xr,description}] \\
\& B_{X;S} \&
\arrow[from=xr,to=xl,bend right=25,"\alpha_X"]
\end{tikzcd}}
with an arrow
\mspage{162}


\scanpage{164}
% Lignes 1--5.
``of commutation'' $\alpha$ which is such only by abuse of language --- it is not an isomorphism \deleted{of \unsure{functors}} but a morphism of functors and nothing more
\[
 \sigma^*\psi^*\deleted{\unsure{u}}\xrightarrow{\alpha_X}\rho^*\varphi^*.
\]
% Lignes 6--11.
In the same way, gluing $D^*$ and $S$ via $\sigma_*$ (i.e. replacing $B_U$ by $B_{D^*}$ in the preceding construction of $\unsure{D_{X;S}|B_{X;S}}$), we find a topos, not multi-Galois in general, denoted $B_{D;S}$. In the \deleted{\unsure{present}} geometric model where we have $X,S$ as above,
% Lignes 12--18.
$B_X$ corresponds to sheaves on $X$ which are \unsure{essentially|set|in} locally constant on $U$ (and on $S$, where they have no merit) i.e. on $U$ arise by inverse image \deleted{of a sheaf} \inserted{under} $U\longrightarrow B_{\Pi_1 U}$ from a sheaf on $B_{\Pi_1 U}$.

% Lignes 19--25.
and in the same way $B_{D;S}$ corresponds to sheaves on $D=\coprod_i\operatorname{Spec}(\mathcal O_i=\text{henselization of }\mathcal O_{X,\unsure{i|s}})$ \deleted{which are} which are locally constant on $D^*=D-S$ (and on $S$, without merit!) but with the dimension hypothesis made -- --- \unsure{in|we} fact
\[
 B_{D;S}\simeq D_{\text{ét}}
\]
($D_i$ has only 2 points, $D_i^*=\{\eta_i\}$ --- ). Thus $B_{D;S}$ \unsure{$(\simeq D_{\text{ét}})$}
\mspage{163}


\scanpage{165}
\mspage{82}
% Lignes 1--2 et diagramme supérieur.
(when we start from $X,S$) fits into a triangle of morphisms of topoi
\diagram{\begin{tikzcd}[ampersand replacement=\&,row sep=2em,column sep=2.5em]
\& B_{D^*} \arrow[dl,equals,""{name=deq,description}] \arrow[dr,"\sigma",""{name=dsig,description}] \& \\
B_{D^*} \arrow[dr,"\varphi_D"'] \& \& B_S \arrow[dl,"\psi_D"] \\
\& B_{D;S} \&
\arrow[from=dsig,to=deq,bend left=22,"\alpha_D"']
\end{tikzcd}}
\margin{\uline{N.B.} in the geometric case, \deleted{\illegible{1}} \unsure{three} topoi of this diagram are étale topoi ($D^*_{\text{ét}},D_{\text{ét}},S_{\text{ét}}$) and the arrows $\varphi_D$ and $\psi_D$ correspond to \unsure{morphisms} \unsure{but not} \illegible{1} of topoi)}
% Lignes 3--4.
$\alpha_D$, like $\alpha_X$ above, is just a plain homomorphism (not an isomorphism in general).

% Lignes 5--14.
\deleted{We thus find a \unsure{topos.}}
When we speak of the canonical morphism from $B_{D^*}$ to $B_{D;S}$ it is \inserted{of} $\varphi_D$ (and not of $\psi_D\sigma$) that we are speaking --- in the geometric case we have $B_{D^*}\simeq D^*_{\text{ét}}$, and $\varphi_D$ corresponds to the inclusion of schemes
\[
 D^*\unsure{\supset|=}D\setminus S\longrightarrow D
\]
\inserted{and $\psi_D:B_S\longrightarrow B_{D;S}$ to the inclusion of schemes $S\longrightarrow D$}
whereas neither $\sigma$ nor $\psi_D\sigma$ correspond in general to morphisms of schemes.

% Lignes 15--16 et diagramme inférieur entièrement barré.
We thus obtain a diagram of topoi
\deleted{\diagram{\begin{tikzcd}[ampersand replacement=\&,row sep=5em,column sep=8em]
\& B_{D^*}
\arrow[dl,"\unsure{\rho_X}"',"\deleted{\unsure{B_{D^*}}}"{description,pos=0.56},"\deleted{\unsure{\varphi_D}}"{pos=0.25}]
\arrow[dl,bend right=38,"\deleted{\unsure{\illegible{1}}}"']
\arrow[dr,"\sigma"{xshift=0.65em,yshift=0.55em}]
\arrow[d] \& \\
B_U \arrow[dr,"\varphi"',""{name=oldl,description}] \& B_D \arrow[d] \& B_S \arrow[l] \arrow[dl,"\psi",""{name=oldr,description}] \\
\& B_X \&
\arrow[from=2-1,to=2-2,phantom,"\text{commutative}"{description,yshift=1.3em}]
\arrow[from=2-2,to=2-3,phantom,"{\begin{gathered}\text{pseudo-}\\[-0.4ex]\text{commutative}\end{gathered}}"{description,xshift=-1.5em,yshift=2.2em}]
\arrow[from=2-2,to=2-3,phantom,"\text{commutative}"{description,xshift=-.8em,yshift=-1.3em}]
\arrow[from=oldr,to=oldl,bend right=20,"\deleted{\unsure{\illegible{1}}}"{pos=.8,xshift=-1.5em,yshift=-.2em}]
\end{tikzcd}}
where in the geometric case}
\mspage{164}


\scanpage{166}
% Diagramme supérieur (**).
\[
 (**)\qquad
 \diagram{\begin{tikzcd}[ampersand replacement=\&, column sep=large,row sep=small]
 \& \& \uline{B}_S\arrow[dl]\arrow[dd,"{\Vert}" description] \\
 \uline{B}_{D^*}\arrow[r]\arrow[dd]\arrow[urr] \& B_{D;S}\arrow[dd] \& \\
 \& \& \uline{B}_S\arrow[dl] \\
 \uline{B}_{\unsure{U}}\arrow[r] \& B_{X;S} \&
 \end{tikzcd}}
\]
\margin{Careful, the triangle is not essentially commutative, but \unsure{we have} pseudo-commutativity\ldots}
% Lignes 1--5.
where the three underlined topoi $\uline{B}_U$, $\uline{B}_S$ and $\uline{B}_{D^*}$ are multi-Galois, and $B_{D;S}$ and $B_{X;S}$ are composite (obtained by gluing two multi-Galois topoi).
\inserted{Unless I am mistaken the two squares are [\uline{not} only essentially$^t$ commutative but also] 2-cartesian (in the 2-category of topoi).}

% Lignes 6--12 et diagramme médian.
Passing to the $B_{\Pi_1}$? of the topoi considered (corresponding to the locally constant objects on these topoi) we find \inserted{a topos} $B_{\deleted{\Pi_1}X}$ \inserted{$=\unsure{B_{\Pi_1 B_{X,S}}|B_{\Pi_1 B_{X;S}}}$} as an amalgamated sum of topoi in the \deleted{$B_{\Pi_1 X}$} diagram
\diagram{\begin{tikzcd}[ampersand replacement=\&,row sep=2em,column sep=2.5em]
\& B_{D^*} \arrow[dl] \arrow[dr] \& \\
B_U \arrow[dr,""{name=pl,description}] \& \& B_S \arrow[dl,no head,""{name=pr,description}] \\
\& B_{\deleted{\unsure{\Pi_1}}X} \&
\arrow[from=pr,to=pl,bend right=25,"\sim"]
\end{tikzcd}}
% Lignes 13--19.
(this square is indeed essentially commutative) and in the same way (replacing $B_U$ by $B_{D^*}$) a $B_{\deleted{\Pi_1}D}$, which is however (via $B_S\longrightarrow B_{\deleted{\Pi_1}D}$) isomorphic \inserted{(rather: \unsure{equivalent|equal})} to $B_S$. The diagram $(**)$ then becomes
\mspage{165}


\scanpage{167}
\mspage{83}
% Diagramme supérieur (***).
\[
 (***)\qquad
 \diagram{\begin{tikzcd}[ampersand replacement=\&, column sep=large,row sep=small]
 \& \& B_S\;(=B_{\Pi_1 S})\arrow[dl,"\sim"']\arrow[dd,"{\Vert}" description] \\
 B_{D^*}\;(=B_{\Pi_1 D^*})\arrow[r]\arrow[dd]\arrow[urr,"\sigma"] \& B_{\deleted{\Pi_1}D}\arrow[dd] \& \\
 \& \& B_S\arrow[dl] \\
 B_U\;(=B_{\Pi_1 U})\arrow[r] \& B_{\deleted{\Pi_1}X} \&
 \end{tikzcd}}
\]
% Lignes 1--4.
where this time the upper triangle is indeed essentially surjective (that is indeed how we defined $\sigma$, in the geometric case!) and we have a morphism from
% Lignes 5--12.
$(**)$ to $(***)$, which \deleted{\unsure{is}} for the arrows which \unsure{are|are} not essentially \inserted{\unsure{(semi-}}identities), $B_{D;S}\longrightarrow B_{\deleted{\Pi_1}D}$ and $B_{X;S}\longrightarrow B_{\deleted{\Pi_1}X}$, have a marked tendency to be ``acyclic'' i.e. to induce isomorphisms on cohomology (we should look at \illegible{1} this point). Finally,
% Lignes 13--15 et diagramme inférieur (****).
in the geometric case, we have \inserted{\unsure{some}} \unsure{analogous|canonical} diagrams of morphisms of étale topoi \deleted{\illegible{1}}
\[
 (****)\qquad
 \diagram{\begin{tikzcd}[ampersand replacement=\&, column sep=large,row sep=small]
 \& \& S_{\text{ét}}\arrow[dl]\arrow[dd,equal] \\
 D^*_{\text{ét}}\arrow[r]\arrow[dd]\arrow[urr] \& D_{\text{ét}}\arrow[dd] \& \\
 \& \& S_{\text{ét}}\arrow[dl] \\
 U_{\text{ét}} \& X_{\text{ét}}\arrow[l] \&
 \end{tikzcd}}
\]
\margin{upper triangle not essentially commutative.}
\mspage{166}


\scanpage{168}
% Lignes 1--2.
where all the arrows except $D^*_{\text{ét}}\longrightarrow S_{\text{ét}}$ are induced by morphisms of topoi.
% Lignes 3--11.
Here $(****)$ maps to the corresponding diagram $(**)$, inducing isomorphisms of topoi for $D^*,D,S$, and for $U_{\text{ét}}\longrightarrow B_U$ and $X_{\text{ét}}\longrightarrow X$ inducing morphisms which have less tendency to be acyclic, but which are \deleted{\unsure{however}} all the same in important cases, recalled in an earlier section\ldots

% Lignes 12--15 et diagramme.
The constructions of $(**)$ and $(***)$ in terms of the starting diagram
\diagram{\begin{tikzcd}[ampersand replacement=\&]
 \& B_{D^*}\arrow[dl]\arrow[dr] \& \\
 B_U \& \& B_S
\end{tikzcd}}
are purely formal, and independent of \unsure{every} hypothesis.
% Lignes 16--22.
The construction of a multi-Galois $B_{\deleted{\Pi_1}X}$, as an amalgamated sum, can be interpreted as the translation (at the level of fundamental groupoids) of an operation of ``plugging holes''.
\deleted{At the level} In the computational context
\mspage{167}


\scanpage{169}
\mspage{84}
% Ligne 1 et diagramme de groupes.
(with a choice of the $\widetilde U,\widetilde U_i$, $\rho_!(\widetilde U_i)\simeq\widetilde U$) we have
\diagram{\begin{tikzcd}[ampersand replacement=\&]
 \& \mathcal E_i\arrow[dl,"\rho_i"']\arrow[dr,"\sigma_i"] \& \\
 \pi \& \& \Gamma_i
\end{tikzcd}},
% Lignes 2--8.
putting $\widetilde X=\text{image of }\widetilde U$ under $i_!$ $(i:B_U\longrightarrow B_{\Pi_1 X})$ $B_{\deleted{\Pi_1}X}$ is described in terms of this $\widetilde U$ as the \deleted{topos of $\Pi$-sets} classifying topos $B_{\pi_X}$, where $\pi_X=\pi_1(X;\widetilde X)$ \deleted{is defined} \inserted{is computed} as the quotient of $\pi$ by the invariant subgroup generated by the $\rho_i(L_i)$, where $L_i=\operatorname{Ker}\sigma_i\subset\mathcal E_i$.

% Lignes 9--16.
If we are given a direct summand $B_{D_0^*}$ of $B_{D^*}$ (corresponding to a subset $I_0$ of $I=\pi_0(B_{D^*})$) we find in the same way an amalgamated sum of $B_U$ and of $B_{S_0}$ under $B_{D_0^*}$, denoted \deleted{$\unsure{B_{\Pi_1(U\cup S_0)}}$} \inserted{$B_{U\cup S_1}\longleftarrow B_U$}, which is visualized as a partial plugging of holes, interpreted at the level of fundamental groupoids. In the geometric case, if we put $I=I_0\amalg I_1$, i.e. $S=S_0\amalg S_1$, we
\mspage{168}


\scanpage{170}
% Lignes 1--2.
can interpret this topos as $B_{\Pi_1 U_1}$, where $U_1=X\setminus\unsure{S_1}$. \deleted{Of course,}

% Lignes 3--6 et premier petit diagramme.
Of course, we have a homomorphism of cocartesian diagrams of topoi relative to
\diagram{\begin{tikzcd}[ampersand replacement=\&]
 \& B_{D_0^*}\arrow[dl]\arrow[dr] \& \\
 B_U \& \& B_{S_0}
\end{tikzcd}}
\deleted{\unsure{from below}} the one relative to $B_{D^*}$ \deleted{\unsure{the situation}} mapping into $B_U$, and \deleted{$\unsure{B_{S_0}}$} $B_S$
% Grand diagramme central, avec surcharges et traits peu lisibles.
\diagram{\begin{tikzcd}[ampersand replacement=\&,row sep=3em,column sep=1.3em]
[B_{D_1^*}] \arrow[dr,dotted] \& \& B_{D_0^*} \arrow[dl] \arrow[dr] \arrow[rrr] \& \& \& B_{D^*} \arrow[dl,phantom,""{name=diagupper,description,pos=0.68,inner sep=0pt},""{name=diaglower,description,pos=0.86,inner sep=0pt}] \arrow[dr] \& \\
\& |[alias=leftU]| B_U \arrow[dr] \arrow[rrr,phantom,""{name=ugapleft,description,pos=0.39},""{name=ugapright,description,pos=0.76},""{name=utip,description,pos=0.88},""{name=uheadstart,description,pos=0.93,inner sep=0pt},""{name=uheadend,description,pos=0.97,inner sep=0pt},""{name=udashleft,description,pos=0.74,yshift=2pt,inner sep=0pt},""{name=udashright,description,pos=0.80,yshift=2pt,inner sep=0pt}] \& \& |[yshift=1.5em]| B_{S_0} \arrow[dl] \arrow[rrr] \& |[alias=rightU]| B_U \arrow[dr] \& \& |[yshift=1.5em]| B_S \arrow[dl] \\
\& \& B_{\deleted{\unsure{\Pi_1}}(\underbrace{U\cup S_0}_{U_1})} \arrow[rrr,no head] \& \& \&
\begin{gathered}
B_{\deleted{\unsure{\Pi_1}}(\underbrace{U\cup S}_{X})}\,\unsure{'}\\[-0.2ex]
\deleted{\unsure{\illegible{1}}\;B_{\unsure{\Pi_1}X}}
\end{gathered} \&
\arrow[from=1-6,to=diagupper,no head,shorten >=0pt]
\arrow[from=diaglower,to=rightU,no head,shorten <=0pt]
\arrow[from=leftU,to=ugapleft]
\arrow[from=ugapright,to=utip,no head]
\arrow[from=uheadstart,to=uheadend,shorten <=0pt,shorten >=0pt]
\arrow[from=udashleft,to=udashright,no head,shorten <=0pt,shorten >=0pt]
\end{tikzcd}}
,
% Lignes 7--9.
and on the other hand we have \inserted{a composite} $B_{D_1^*}\longrightarrow B_U\longrightarrow B_{\deleted{\Pi_1}U_1}$ which with $\unsure{B_{D_1^*}}\longrightarrow B_{S_1}$ \deleted{can be used} \inserted{gives a diagram} of topoi
\diagram{\begin{tikzcd}[ampersand replacement=\&]
 \& B_{D_1^*}\arrow[dl]\arrow[dr] \& \\
 B_{\deleted{\Pi_1}U_1} \& \& B_{S_1}
\end{tikzcd}}
% Lignes 10--15 (phrase continuée hors du lot).
of the same type as at the beginning, which we can use to construct another \deleted{\unsure{thus|in}} amalgamated sum. And it is obvious that this one is canonically equivalent to $B_{\deleted{\Pi_1}X}$, the amalgamated sum corres-
\mspage{169}
